% s-364-p364-m002; original PDF 364, 365, 366
\Aphorism{s-364-p364-m002}{408}
Der naive Realismus, der den scheinbaren, empirischen Weltverlauf für eine unmittelbar getreue Abbildung des wirklichen, transcendenten hält — was für eine ungeheure Zumuthung stellt er eigentlich an die Weisheit und Harmonie dieses transcendenten Verlaufs! Man versuche im Ernst, die Complication und wunderbar verzweigte Structur unseres Kosmos nachzudenken, im Kleinen und Grossen, im Ganzen und Einzelnen, von den Vorgängen in der Infusorienwelt eines Wassertropfens bis hinauf zu den Kreisläufen der Sonnen- und Milchstrassensysteme. Mit welch umfassendem und gegliedertem Inhalt ist eine einzige Secunde empirischen Geschehens erfüllt: in dieser Secunde beschreibt die Sonne ein Stück ihrer Bahn im Astralgebäude, die Erde und sämmtliche Planeten ein Stück ihrer Bahnen um die Sonne, Monde, Kometen und unzählige Meteorsplitter ebenfalls ihre vorgeschriebenen Wegfragmente, sammt und sonders unter Innehaltung bestimmter Gesetze und Eins auf das Andere Rücksicht nehmend und Fernkraft ausübend. In dieser selben Secunde fällt auf Erden jeder Regentropfen, wie er fallen soll, wächst jede Blume um ein Bestimmtes, vollzieht jedes Insect nothwendige und determinirte Bewegungen, tauchen in jedem Menschengehirn Vorstellungen auf, die das genaue Resultat von Vergangenheit und Umgebung sind — und dieser ganze vieltönige vielfarbige vielgestaltige Mechanismus soll sich in \Sperr{Wirklichkeit} so abspielen, mit minutiöser Beachtung aller der Millionen Nothwendigkeiten und „Naturgesetze“, die hier bunt durcheinander commandiren? Welch eine Virtuosität, welch ein Aufmerken, Blicken, Lenken nach allen Seiten wird da vorausgesetzt, welch ein Seiltanz wäre nöthig, um hier nicht zu fallen! welch ein Dirigent müsste dies unendlich vielstimmige Orchester zusammenhalten! Und schliesslich mündet der naive Realismus nie sehr weit von dem Ziel, eben den Dirigenten zu \Sperr{beweisen}, den Gott hinter der Natur zu beweisen; er ist das physico-theologische Argument oder die Vorbereitung dazu. Dagegen lehrt der Idealismus, dass der zweckmässige, causale, reichorganisirte Verlauf der empirischen Wirklichkeit gar nichts über den transcendenten besagt, dass ein „an sich“ willkürliches Geschehen ausreicht, uns als subjectiven Effect ein determinirtes Geschehen zu garantiren, dass auch der vollkommenste Kosmos nur als Durchsiebungsproduct aus dem wüstesten Chaos hervorgeht. Das Orchester also bedarf keines Dirigenten, es bedarf nicht einmal der Noten oder reingestimmter Instrumente, es kann ungeleitet, disharmonisch durcheinander toben: wir, die Hörer, haben jene wunderbare Vorrichtung in den. Ohren, vermöge deren wir stets eine rhythmisch klare, rein harmonisirte Melodie vernehmen. Wie das? wir hören überhaupt nur, wenn jene \Sperr{kosmische} Ausnahme-Musik ertönt, alle übrige Zeit sind wir nicht allein taub, sondern existenzlos. Sobald wir nun uns selbst als existirend vorfinden — und wann könnten wir uns existenzlos antreffen? — finden wir mit uns auch den Kosmos vor, der unsere Existenzbedingung ist: d. h. wir glauben uns beständig von Kosmos, nie von Chaos umgeben. Diese Vorstellung, dass eine kosmisch geordnete Welt für uns geradezu Bewusstseinsbedingung ist, mag im Einzelnen schwer zu belegen sein; als Ganzes ist sie der einzige philosophische Gedanke, mit dem wir um die Thorheit herumkommen, vom transcendenten Weltverlauf die Specialitäten und Gesetzlichkeiten unseres empirischen zu verlangen. Es hat keine philosophische Schwierigkeit, das Bewusstsein an Bedingungen zu knüpfen und höheres Bewusstsein an complicirtere Bedingungen; aber es ist nicht nur schwierig, sondern geradezu absurd, ein „Ding an sich“, d. h. eine von unserem Bewusstsein unabhängige Realität anzusetzen und hinterher diese Unabhängigkeit zu einem blossen Parallelismus mit unserem Bewusstsein einzuschränken.
\par

